% Physics formula booklet: every formula of the IBE3 Physics for Engineers notes
% Electricity Part 1 (Boxes 1-24) and Thermodynamics Part 1 (Boxes 1-18), with symbols,
% units and conditions, plus a marked part on energy transfers and heat engines.
% Entries marked with a dagger (\fe*) are not printed in the notes.
\documentclass[10pt,a4paper]{article}
\usepackage{formulabook}
\colorlet{theme}{navy}
\renewcommand{\booktitle}{Physics for Engineers \textperiodcentered{} Formula booklet}
\newcommand{\ETh}{E_{\mathrm{Th}}}
\newcommand{\RTh}{R_{\mathrm{Th}}}
\newcommand{\INo}{I_{\mathrm{No}}}
\newcommand{\RNo}{R_{\mathrm{No}}}
\newcommand{\Req}{R_{\mathrm{eq}}}
\newcommand{\ist}{i^{\star}}
\newcommand{\ust}{u^{\star}}
\newcommand{\dd}{\mathrm d}
\newcommand{\pd}[3]{\left(\frac{\partial #1}{\partial #2}\right)_{\!#3}}

\begin{document}
\thispagestyle{first}

\begin{tcolorbox}[enhanced, colback=theme, colframe=theme, coltext=white, arc=2mm,
    boxrule=0pt, left=10pt, right=10pt, top=9pt, bottom=9pt]
  {\LARGE\bfseries Physics for Engineers}\\[3pt]
  {\Large Formula booklet}\\[6pt]
  {\small Electricity Part 1: foundations of DC circuit theory \textperiodcentered{}
  Thermodynamics Part 1: description of a system \textperiodcentered{} heat engines}
\end{tcolorbox}

\begin{keybox}[title={How to read this booklet}]
\small
Each row gives a quantity and where it is in the notes (Box number, equation or page), the
formula, then the symbols with their units and when the formula applies. Page numbers refer to
the lecture notes, whose blank boxes are filled in here.
\begin{itemize}
  \item \dagnote\ marks a result that is \textbf{not printed in the notes}: a direct
    consequence, a standard result used in the exercises, or the material of Part~III. Use
    it when it helps, but quote the notes' form when there is one.
  \item Part~III (work, heat, heat engines) goes beyond Part~1 of the Thermodynamics notes.
    It includes the heat taken in by an isothermal expansion,
    $Q = nRT_H\ln(V_{\max}/V_{\min})$.
\end{itemize}
\end{keybox}

\begin{minipage}[t]{0.47\textwidth}
\fbsub{Contents}
\fbcontents
\end{minipage}\hfill
\begin{minipage}[t]{0.49\textwidth}
\fbsub{Constants}
{\small
\renewcommand{\arraystretch}{1.12}
\begin{tabular}{@{}l@{\hspace{6pt}}l@{}}
Elementary charge & $e = 1.6\times10^{-19}$~C\\
Avogadro constant & $N_A = 6.022\times10^{23}$~mol$^{-1}$\\
Ideal gas constant & $R = 8.314$~J\,mol$^{-1}$\,K$^{-1}$ ($\approx 8.3$)\\
Boltzmann constant & $k_B = R/N_A = 1.38\times10^{-23}$~J\,K$^{-1}$\dagnote\\
Gravity & $g = 9.81$~m\,s$^{-2}$\\
At 300~K & $RT \approx 2.5\times10^{3}$~J\,mol$^{-1}$\\
\end{tabular}}

\fbsub{Conversions}
{\small
\renewcommand{\arraystretch}{1.12}
\begin{tabular}{@{}l@{}}
$T\,(\mathrm K) = \theta\,(^\circ\mathrm C) + 273.15$\\
$1~\mathrm{bar} = 10^{5}$~Pa;\quad $1~\mathrm{atm} = 1.013\times10^{5}$~Pa\\
$1~\mathrm L = 1~\mathrm{dm}^3 = 10^{-3}$~m$^3$;\quad $1~\mathrm{mL} = 1~\mathrm{cm}^3 = 10^{-6}$~m$^3$\\
$1~\mathrm{mm}^2 = 10^{-6}$~m$^2$;\quad $1~\mathrm{cm}^2 = 10^{-4}$~m$^2$\\
$1~\mathrm{W\,h} = 3600$~J;\quad $1~\mathrm{kW\,h} = 3.6\times10^{6}$~J\\
$1~\mathrm{A\,h} = 3600$~C;\quad $1~\mathrm{k\Omega}\times1~\mathrm{mA} = 1$~V\\
\end{tabular}}

\fbsub{SI prefixes}
{\small
\renewcommand{\arraystretch}{1.12}
\begin{tabular}{@{}lll@{}}
G (giga) $10^{9}$ & M (mega) $10^{6}$ & k (kilo) $10^{3}$\\
m (milli) $10^{-3}$ & \textmu{} (micro) $10^{-6}$ & n (nano) $10^{-9}$\\
p (pico) $10^{-12}$ & & \\
\end{tabular}}\\[3pt]
{\footnotesize Exact symbols: k (kilo) $\neq$ K (kelvin), m (milli) $\neq$ M (mega).}
\end{minipage}

% =====================================================================================
\partheading{Part I}{Electricity: foundations of DC circuit theory}
\fbnote{Notes Elec-S1 Part~1 (38 pages). Arrows: the arrow of $u_{AB}$ points towards $A$.
  Receiver convention: $u$ and $i$ arrows opposite; generator convention: same direction.}

\fbsection{1.A}{Electric voltage (pp.~2--7)}
\fe[Box 1]{Electric voltage}{u_{AB} = V_A - V_B}
  {Volt (V). Only potential differences can be measured; potential and voltage have the
   same unit.}
\fe[eq.~(1)]{Sign of a voltage}{u_{AB} = -u_{BA}}
  {$u_{AB} > 0 \Leftrightarrow V_A > V_B$.}
\fe*{Adding voltages}{u_{AC} = u_{AB} + u_{BC}}
  {Potentials subtract like heights; the basis of the mesh law.}
\fe[p.~2]{Ground}{V_{\text{ground}} = 0}
  {Reference of potentials: the wire on the $-$ terminal of the source. ``Earth'' if it is
   connected to the Earth.}
\fe[Box 2]{Ideal voltage source}{u = e\quad\text{whatever } i}
  {$e$: electromotive force, in volts (not newtons). A real source has an internal
   resistance: $u$ drops as $i$ grows (see the Thévenin model, 3.B).}
\fe[p.~7]{Voltmeter}{R_V \to \infty}
  {In parallel, COM on the $-$ side (otherwise the sign flips); $R_V = 1$ to 10~M$\Omega$.}
\fbnote{\textbf{Orders of magnitude} (Tab.~1): neuron $\sim$100~mV; AA cell 1.5~V; Li-ion
  3.6--4.2~V; mains 230~V$_{\text{rms}}$; high-speed train line 25~kV; extra-high-voltage line
  400~kV.}

\fbsection{1.B}{Electric current (pp.~7--11)}
\fe[Box 3]{Electric current}{i = \frac{\dd q}{\dd t}}
  {Ampere (A), an SI base unit: $1~\mathrm A = 1~\mathrm{C\,s^{-1}}$. Same current all
   along a series path. Conventional direction: that of positive charges.}
\fe*{Charge transported}{q = \int_{t_1}^{t_2} i\,\dd t\qquad\text{DC: } q = i\,\Delta t}
  {Coulomb (C).}
\fe*{Number of electrons}{N = \frac{|q|}{e} = \frac{i\,\Delta t}{e}}
  {$e = 1.6\times10^{-19}$~C.}
\fe[Box 4, eq.~(2)]{Microscopic expression}{i = n\,q\,v\,S}
  {$n$ carrier density (m$^{-3}$), $q$ carrier charge (C), $v$ drift velocity
   (m\,s$^{-1}$), $S$ cross-section (m$^2$). Electrons: $i = -n e v S$.}
\fe*{Drift velocity}{v = \frac{i}{n\,e\,S}}
  {Copper, 1~mm$^2$, 1~A: $v \approx 7\times10^{-5}$~m\,s$^{-1}$. Not the speed of the
   signal.}
\fe*{Current density}{j = \frac{i}{S} = n\,|q|\,v}{A\,m$^{-2}$.}
\fe[p.~11]{Ammeter}{R_A \to 0}
  {In series (open the branch to insert it); $R_A < 1~\Omega$.}

\fbsection{1.C}{Circuits and two-terminal components (pp.~11--14)}
\fe[Box 5]{Circuit vocabulary}{\begin{gathered}\text{node: junction of} \ge 3 \text{ conductors}\\
  \text{branch: between two nodes}\end{gathered}}
  {Loop: any closed path. Mesh: a loop containing no other loop. Points joined only by
   wire form one (electrically distinct) node.}
\fe[3.A]{Counting equations}{N - 1 \text{ node laws},\qquad M = B - N + 1 \text{ meshes}}
  {$N$ distinct nodes, $B$ branches: $B$ equations for the $B$ branch currents.}
\fe[Box 6]{Series and parallel}{\begin{gathered}\text{series: same } i,\quad u = \textstyle\sum_k u_k\\
  \text{parallel: same } u,\quad i = \textstyle\sum_k i_k\end{gathered}}
  {Same current in series (node law), voltages add (mesh law). Same voltage in parallel
   (mesh law), currents add at the nodes (node law).}
\fe[p.~13]{Direction of the current}{\text{outside the source: from } \oplus \text{ to } \ominus}
  {Electrons flow the other way ($i = -n e v S$).}
\fe[Box 7]{Electrical power}{p = u\,i}
  {Watt: $1~\mathrm W = 1~\mathrm V\times1~\mathrm A$. Receiver convention: $p > 0$ is
   received; generator convention: $p > 0$ is supplied.}
\fe[eqs.~(4), (5)]{Electrical energy}{E = \int_{t_1}^{t_2} u\,i\,\dd t\qquad\text{DC: } E = u_0\,i_0\,\Delta t}
  {Joule (J); $1~\mathrm{W\,h} = 3600$~J.}

\fbsection{2.A}{Characteristic curve (pp.~14--19)}
\fe[Box 8]{Characteristic}{i = f(u)\quad\text{or}\quad u = g(i)}
  {Pairs $(u_k, i_k)$ measured in many states. Depends on the convention: receiver for a
   passive component, generator for an active one.}
\fe[Box 9]{Resistor}{i = \frac{u}{R}\text{: a line through } O}
  {Slope of $i = f(u)$: $G = 1/R$. Slope of $u = g(i)$: $R$. The signature of an ohmic
   conductor.}
\fe[Box 10]{Quick (naive) estimate of $R$}{R = \frac{1}{N}\sum_{k=1}^{N}\frac{u_k}{i_k}}
  {Mean of the ratios $r_k = u_k/i_k$; treats the noisy low-current points like the others.}
\fe[Box 11]{Estimate by linear regression}{\hat R = \frac{\sum_k u_k\,i_k}{\sum_k i_k^2}}
  {Minimises $\phi(R) = \sum_k (u_k - R\,i_k)^2$: best line $u = \hat R\,i$ through $O$.}
\fe[p.~18]{Standard uncertainty on $\hat R$}{u(\hat R) = \sqrt{\frac{\frac{1}{N-1}\sum_k
  (u_k - \hat R\,i_k)^2}{\sum_k i_k^2}}}
  {$u_k - \hat R\,i_k$: residuals (Tab.~4).}
\fe[Figs.~15, 19]{Other characteristics}{\text{not lines through } O}
  {Ideal voltage source: vertical line $u = E$. Ideal current source: horizontal line
   $i = I_0$. Lamp: curved. Diode, Zener diode ($-U_Z$).}

\fbsection{2.B}{Ohm's law and Joule's law (pp.~19--21)}
\fe[Box 12]{Ohm's law}{u = R\,i}
  {Receiver convention ($u = -R\,i$ in generator convention). $R$ in ohms ($\Omega$).}
\fe[eq.~(6)]{Conductance}{G = \frac{1}{R},\qquad i = G\,u}{Siemens: $1~\mathrm S = 1~\Omega^{-1}$.}
\fe[Box 13, eq.~(7)]{Joule's law}{p = u\,i = R\,i^2 = \frac{u^2}{R}}
  {Power turned into heat by $R$ (W). It is $R\,i^2$, never $i\,R^2$.}
\fe*{Heat released}{E_J = R\,i^2\,\Delta t}{Constant current $i$ for a time $\Delta t$.}
\fbnote{\textbf{Typical resistances} (Tab.~6): lab lead, a few tens of m$\Omega$; ammeter,
  10~m$\Omega$ to a few $\Omega$; lab resistor, 1~$\Omega$ to 1~M$\Omega$; dry skin, a few
  k$\Omega$ to tens of k$\Omega$; voltmeter input, 1 to 10~M$\Omega$.}

\fbsection{2.C}{Pouillet's law (pp.~21--22)}
\fe[Box 14]{Pouillet's law}{R = \rho\,\frac{L}{S}}
  {$\rho$ resistivity ($\Omega$\,m), $L$ length (m), $S$ cross-section (m$^2$). Circular
   wire: $S = \pi d^2/4$.}
\fe[Box 15]{Conductivity}{\sigma = \frac{1}{\rho},\qquad R = \frac{L}{\sigma\,S}}
  {S\,m$^{-1}$. Good conductor: small $\rho$, large $\sigma$.}
\fe*{Scaling}{R \propto \frac{L}{d^2}}{Doubling the diameter divides $R$ by 4.}
\par\smallskip
\begin{center}
\small
\renewcommand{\arraystretch}{1.1}
\begin{tabular}{@{}lccccc@{}}
\toprule
Material (Tab.~7) & Silver & Copper & Aluminium & Iron & Graphite\\
\midrule
$\rho$ ($\Omega$\,m) & $1.6\times10^{-8}$ & $1.7\times10^{-8}$ & $2.8\times10^{-8}$ &
  $1.0\times10^{-7}$ & $\sim3\times10^{-5}$\\
$\sigma$ (S\,m$^{-1}$) & $6.3\times10^{7}$ & $5.9\times10^{7}$ & $3.6\times10^{7}$ &
  $1.0\times10^{7}$ & $\sim3\times10^{4}$\\
\bottomrule
\end{tabular}\\[2pt]
{\footnotesize For most metals $\rho$ increases with temperature. Glass: $\rho$ very large.}
\end{center}

\fbsection{3.A}{Kirchhoff's laws and their consequences (pp.~22--29)}
\fe[Box 16]{Node law}{\sum i_{\text{in}} = \sum i_{\text{out}}}
  {Conservation of charge: nothing piles up at a node.}
\fe[Box 16]{Mesh law}{\sum_{\text{mesh}} \varepsilon_k\,u_k = 0}
  {Conservation of energy. $\varepsilon_k = +1$ if the arrow of $u_k$ points along the
   direction of travel, $-1$ otherwise.}
\fe[Box 17, eq.~(8)]{Resistors in series}{\Req = R_1 + R_2 + \dots + R_n}
  {Larger than the largest $R_k$; $n$ identical: $nR$.}
\fe[Box 18, eq.~(9)]{Resistors in parallel}{\frac{1}{\Req} = \frac{1}{R_1} + \frac{1}{R_2} + \dots + \frac{1}{R_n}}
  {$G_{\text{eq}} = \sum G_k$. Smaller than the smallest $R_k$; $n$ identical: $R/n$.}
\fe[eq.~(10)]{Two in parallel}{\Req = \frac{R_1 R_2}{R_1 + R_2}}{Product over sum.}
\fe[Box 19, eq.~(12)]{Voltage divider}{u_1 = \frac{R_1}{R_1 + R_2}\,u,\qquad u_2 = \frac{R_2}{R_1 + R_2}\,u}
  {Same current in $R_1$ and $R_2$ (nothing drawn from the midpoint). $u_1/u_2 = R_1/R_2$.
   With a load $R_L$ across $R_2$: use $R_2 \parallel R_L$.}
\fe*{Current divider}{i_1 = \frac{R_2}{R_1 + R_2}\,i = \frac{G_1}{G_1 + G_2}\,i}
  {Two resistors in parallel sharing $i$ (crossed indices).}
\fe[Box 20, eq.~(17)]{Millman's theorem}{V_O = \frac{\sum_k V_k/R_k}{\sum_k 1/R_k}}
  {Node $O$ joined only through resistors $R_k$ to potentials $V_k$. Source $E_k$ in the
   branch: $V_k = \pm E_k$; branch to ground: $V_k = 0$ but $1/R_k$ stays below.}

\fbsection{3.B}{Thévenin's and Norton's theorems (pp.~29--33)}
\fe[Fig.~25]{Thévenin model}{u = \ETh - \RTh\,i}
  {Ideal source $\ETh$ in series with $\RTh$. Generator convention: $i$ leaves by $A$,
   $u = V_A - V_B$.}
\fe[Fig.~26]{Norton model}{i = \INo - \frac{u}{\RNo}}{Ideal current source $\INo$ in parallel with $\RNo$.}
\fe[Box 21, eq.~(18)]{Equivalence of the models}{\RTh = \RNo \equiv R,\qquad \ETh = R\,\INo}
  {$\ETh$: open-circuit voltage ($i = 0$). $\INo$: short-circuit current ($u = 0$).}
\fe[Box 22]{Thévenin's theorem}{\ETh = u_{AB}\big|_{i = 0},\qquad \RTh = R_{AB}\big|_{\text{sources off}}}
  {Valid for ideal sources and ohmic conductors only (no lamp, motor, diode).}
\fe[Box 23]{Norton's theorem}{\INo = i_{A\to B}\big|_{u = 0},\qquad \RNo = \RTh}
  {$\INo$: current in a wire joining $A$ and $B$, sources on.}
\fe[p.~32]{Switching a source off}{\begin{gathered}\text{voltage source} \to \text{wire}\ (u = 0)\\
  \text{current source} \to \text{open circuit}\ (i = 0)\end{gathered}}
  {The resistors stay in place.}
\fe*{Source transformation}{E \text{ in series with } R \;\Leftrightarrow\; \frac{E}{R} \text{ in parallel with } R}
  {Thévenin models in series add $E$ and $R$; Norton models in parallel add $I$ and $G$.}
\fe*{Superposition}{x = x_1 + x_2 + \dots}
  {Linear circuit: the response to several sources is the sum of the responses to each
   source alone, the others switched off.}
\fe[p.~33]{Thévenin model of a divider}{\RTh = \frac{R_1R_2}{R_1 + R_2},\qquad \ETh = \frac{R_2}{R_1 + R_2}\,E}
  {$E$ in series with $R_1$, $R_2$ across $A$--$B$. Then $\INo = E/R_1$.}

\fbsection{3.C}{Operating point (pp.~34--38)}
\fe[Figs.~28, 29]{Characteristics}{\text{source: } u = \ETh - \RTh\,i\qquad\text{load: } u = R_c\,i}
  {Operating point $(\ust, \ist)$: the unique pair on both; graphically their intersection.}
\fe[Box 24, eq.~(21)]{Operating point}{\ist = \frac{\ETh}{\RTh + R_c},\qquad \ust = \frac{R_c}{\RTh + R_c}\,\ETh}
  {$R_c \to 0$: $\ust \to 0$, $\ist \to \ETh/\RTh$. $R_c \to \infty$: $\ist \to 0$,
   $\ust \to \ETh$.}
\fe[eqs.~(19), (20)]{Power received by the load}{P_c = \ust\,\ist = R_c\,{\ist}^2 = \frac{R_c\,\ETh^2}{(\RTh + R_c)^2}}
  {Equal to the power delivered by the real source at $A$--$B$.}
\fe[eqs.~(22), (23)]{Power balance}{\ETh\,\ist = \RTh\,{\ist}^2 + R_c\,{\ist}^2}
  {Produced by the ideal source $=$ lost in $\RTh$ $+$ received by the load.}
\fe[eq.~(24)]{Efficiency}{\eta = \frac{P_c}{\ETh\,\ist} = \frac{R_c}{\RTh + R_c}}
  {0 on short circuit, close to 1 for $R_c \gg \RTh$, $\frac12$ at $R_c = \RTh$.}
\fe[eqs.~(25), (26)]{Maximum power (matching)}{R_c = \RTh \;\Longrightarrow\; P_c^{\max} = \frac{\ETh^2}{4\RTh}}
  {$\dd P_c/\dd R_c = \ETh^2(\RTh - R_c)/(\RTh + R_c)^3$. Efficiency only 50\,\% at
   matching: maximum power $\neq$ maximum efficiency.}
\begin{trapbox}[title={Misprints and slips in the Electricity notes}]
\small
\begin{itemize}
  \item Eq.~(17): the denominator of Millman's formula is $1/R_1 + 1/R_2 + \dots + 1/R_n$.
  \item Eqs.~(15)--(16): the left-hand side is $u_2$ (the voltage across $R_2$), not $u_1$;
    ``$0.9901\,E$'' means $0.9901\,u$.
  \item Eq.~(24): $P_{\mathrm{Th}}$ there means $\ETh\ist$ (produced by the ideal source),
    not the delivered power of eq.~(19).
  \item Tabs.~3--4: the five ratios $u_k/i_k$ average 50.04~$\Omega$ (printed 50.17~$\Omega$);
    the regression gives $\hat R = 49.99~\Omega \approx 50.0~\Omega$; with these residuals the
    uncertainty formula gives $u(\hat R) = 0.08~\Omega$ (printed 1.001~$\Omega$).
\end{itemize}
\end{trapbox}

% =====================================================================================
\partheading{Part II}{Thermodynamics: description of a system}
\fbnote{Notes Thermo-S1 Part~1 (23 pages). Closed systems only: $n$ is constant.}

\fbsection{1.A}{System and state variables (pp.~2--5)}
\fe[Boxes 1--3]{Kinds of system}{\begin{gathered}\text{closed: no matter exchanged}\\
  \text{isolated: neither matter nor energy}\end{gathered}}
  {Open otherwise. Isolated $\Rightarrow$ closed, but not the converse.}
\fe[Tab.~1]{State variables of a gas}{p\ (\mathrm{Pa}),\quad V\ (\mathrm m^3),\quad T\ (\mathrm K),\quad n\ (\mathrm{mol})}
  {Dimensions: $\mathrm M\,\mathrm L^{-1}\mathrm T^{-2}$, $\mathrm L^3$, $\Theta$,
   $\mathrm N$. Defined only at equilibrium (Box 4).}
\fe[Box 5, eq.~(1)]{Extensive quantity}{X(2\times\text{system}) = 2\,X}
  {Proportional to the size: $m$, $V$, $n$, $E$.}
\fe[Box 6]{Intensive quantity}{X(2\times\text{system}) = X}
  {$p$, $T$, constants such as $R$. Ratio of two extensive quantities: intensive ($V/n$).}
\fe[Box 7]{Homogeneity}{\text{intensive} = \text{intensive},\qquad \text{extensive} = \text{extensive}}
  {A check besides dimensions: $T = pV/(nR)$ passes, $T = nR$ fails.}

\fbsub{Composition quantities (Tab.~2)}
\fe[Tab.~2]{Mole fraction}{x_i = \frac{n_i}{n_{\text{tot}}},\qquad \sum_i x_i = 1}{No unit.}
\fe[Tab.~2]{Molar density}{C = \frac{n}{V}}{mol\,m$^{-3}$.}
\fe[Tab.~2]{Molar volume}{V_m = \frac{V}{n}}{m$^3$\,mol$^{-1}$ (inverse of $C$).}
\fe[Tab.~2]{Mass density}{\rho = \frac{m}{V}}{kg\,m$^{-3}$ (also written $\mu$).}
\fe[Tab.~2]{Molar mass}{M = \frac{m}{n}}{g\,mol$^{-1}$.}
\fe[Tab.~2]{Mass concentration}{c_m = \frac{m_{\text{solute}}}{V_{\text{solution}}}}{g\,L$^{-1}$ (also $\gamma$).}
\fe[Tab.~2]{Molar concentration}{c = \frac{n_{\text{solute}}}{V_{\text{solution}}}}
  {mol\,L$^{-1}$ (also [X]). $c_m = c\,M$.}
\fe*{Number of entities}{N = n\,N_A}{$N_A = 6.022\times10^{23}$~mol$^{-1}$.}

\fbsection{1.B}{Postulates of thermodynamics (pp.~5--7)}
\fe[eqs.~(2), (3)]{Thermodynamic equilibrium}{T_{\text{system}} = T_{\text{ext}},\qquad p_{\text{system}} = p_{\text{ext}}}
  {Thermal and mechanical equilibrium: uniform and steady state variables.}
\fe[Box 8]{Postulate}{\begin{gathered}\text{a few state variables}\\ \text{describe an equilibrium state}\end{gathered}}
  {For a closed system: $p$, $V$, $T$.}
\fe*{Pressure}{p = \frac{F}{S}}{Pa $=$ N\,m$^{-2}$; force $F$ normal to the area $S$.}
\fe[Fig.~5]{Piston carrying a mass}{p = p_0 + \frac{m\,g}{S}}
  {Mechanical equilibrium of a vertical piston of mass $m$ and area $S$ under the outside
   pressure $p_0$.}
\fe*{Hydrostatic pressure}{p = p_0 + \rho\,g\,h}{Depth $h$ in a liquid of density $\rho$.}

\fbsection{1.C}{Equations of state and state functions (pp.~7--10)}
\fe[Box 9, eq.~(4)]{Equation of state}{\psi(p, V, T) = 0 \;\Longleftrightarrow\; V = V(T, p)}
  {Valid only at equilibrium; characterises the material.}
\fe[eq.~(5)]{Differential of a state function}{\dd V = \pd{V}{T}{p}\dd T + \pd{V}{p}{T}\dd p}
  {Each partial derivative is the slope of an isobar $V(T)$ or an isotherm $V(p)$ (Fig.~7).}
\fe[eq.~(11)]{Cyclic rule}{\pd{p}{T}{V}\pd{T}{V}{p}\pd{V}{p}{T} = -1}{For any relation $f(p, V, T) = 0$.}
\fe*{Inverse derivative}{\pd{V}{p}{T} = \frac{1}{(\partial p/\partial V)_T}}{Same variable held fixed.}

\fbsection{2.A}{Condensed phases: $\chi_T$ and $\alpha$ (pp.~11--13)}
\fe[Box 10, eq.~(7)]{Isothermal compressibility}{\chi_T = -\frac{1}{V}\pd{V}{p}{T} = \frac{1}{\rho}\pd{\rho}{p}{T}}
  {Pa$^{-1}$, positive. Condensed phases: $\chi_T \approx 0$.}
\fe[Box 11]{Isobaric expansivity}{\alpha = \frac{1}{V}\pd{V}{T}{p}}
  {K$^{-1}$; nearly constant over a limited range of $T$. Also
   $\alpha = -\frac1\rho\big(\frac{\partial\rho}{\partial T}\big)_p$.}
\fe*{Relative variations}{\frac{\dd V}{V} = \alpha\,\dd T - \chi_T\,\dd p}
  {Small changes: $\frac{\Delta V}{V} \approx \alpha\,\Delta T - \chi_T\,\Delta p$, and
   $\frac{\Delta\rho}{\rho} = -\frac{\Delta V}{V}$.}
\fe[Box 13]{Equation of state of a condensed phase}{\begin{gathered}V = V_0\,e^{\alpha(T - T_0) - \chi_T(p - p_0)}\\
  \approx V_0\left[1 + \alpha(T - T_0) - \chi_T(p - p_0)\right]\end{gathered}}
  {$\alpha$, $\chi_T$ constant; integrate $\dd V/V$. Limit $\alpha = \chi_T = 0$: $V = V_0$.}
\fe*{Heating at constant volume}{\pd{p}{T}{V} = \frac{\alpha}{\chi_T}\qquad \Delta p \approx \frac{\alpha}{\chi_T}\,\Delta T}
  {Liquid sealed in a rigid container ($\dd V = 0$).}
\fbnote{\textbf{Orders of magnitude}\dagnote: water $\alpha \approx 2\times10^{-4}$~K$^{-1}$,
  $\chi_T \approx 5\times10^{-10}$~Pa$^{-1}$; ideal gas at 1~bar and 300~K
  $\chi_T = 10^{-5}$~Pa$^{-1}$, $\alpha = 3.3\times10^{-3}$~K$^{-1}$.}

\fbsection{2.B}{Gases: $\beta$ (pp.~13--14)}
\fe[Box 12]{Piezothermic coefficient}{\beta = \frac{1}{p}\pd{p}{T}{V}}
  {K$^{-1}$: relative rise of the pressure with $T$ at constant volume.}

\fbsection{2.C}{Principle and method (pp.~14--15)}
\fe[eq.~(11)]{Link between the three}{\beta = \frac{\alpha}{p\,\chi_T}\quad\Longleftrightarrow\quad \alpha = p\,\beta\,\chi_T}
  {From the cyclic rule applied to $f(T, V, p) = 0$.}
\fe[Fig.~8]{Method}{\text{equation of state} \;\xrightleftharpoons[\text{integrate}]{\text{differentiate}}\; \alpha,\ \chi_T,\ \beta}
  {Coefficients measured by experiment give the equation of state, and conversely.}

\fbsection{3.A}{Low-density gases: the historical laws (pp.~15--17)}
\fe[eq.~(12)]{Boyle--Mariotte (1662)}{pV = \text{const}\qquad(n, T \text{ fixed})}
  {$p_1V_1 = p_2V_2$. Gives $\chi_T = 1/p$ (Box 14).}
\fe[eq.~(13)]{Charles (1787)}{\frac{V}{T} = \text{const}\qquad(n, p \text{ fixed})}
  {$V_1/T_1 = V_2/T_2$. Gives $\alpha = 1/T$ (Box 15).}
\fe[eq.~(14)]{Gay-Lussac (1802)}{\frac{p}{T} = \text{const}\qquad(n, V \text{ fixed})}
  {$p_1/T_1 = p_2/T_2$. Gives $\beta = 1/T$ (Box 16).}
\fe[p.~15]{Avogadro (1811)}{\frac{V}{n} = \text{const}\qquad(p, T \text{ fixed})}
  {Same $p$, $T$ and $V$ $\Rightarrow$ same amount of gas.}

\fbsection{3.B}{The ideal gas (pp.~17--20)}
\fe[Box 17]{Ideal gas law (Clapeyron)}{pV = nRT}
  {$p$ in Pa, $V$ in m$^3$, $T$ in K, $R = 8.314$~J\,mol$^{-1}$\,K$^{-1}$. Limit of every gas
   when $n/V \to 0$ (no interactions).}
\fe[Box 18]{Derivation}{\frac{pV}{T} = \text{const} = nR}
  {Combine Boyle--Mariotte, Charles and Gay-Lussac; Avogadro gives the factor $n$.}
\fe*{Two states, $n$ fixed}{\frac{p_1V_1}{T_1} = \frac{p_2V_2}{T_2}}
  {Identify first which variables are fixed (table below).}
\fe*{Other forms}{pV_m = RT,\qquad \rho = \frac{pM}{RT},\qquad C = \frac{p}{RT}}
  {$M$ in kg\,mol$^{-1}$ for $\rho$ in kg\,m$^{-3}$.}
\fe*{Per molecule}{pV = N\,k_B\,T}{$k_B = R/N_A = 1.38\times10^{-23}$~J\,K$^{-1}$.}
\fe[Boxes 14--16]{Coefficients of the ideal gas}{\chi_T = \frac1p,\qquad \alpha = \frac1T,\qquad \beta = \frac1T}
  {Check: $\alpha/(p\,\chi_T) = 1/T = \beta$.}
\fe*{Molar volume}{V_m = \frac{RT}{p}}
  {22.4~L\,mol$^{-1}$ at 0~$^\circ$C and 1~atm; 24.8~L\,mol$^{-1}$ at 25~$^\circ$C and 1~bar.}
\fe*{Mixture (Dalton)}{p_i = x_i\,p,\qquad p = \sum_i p_i}
  {Partial pressures. Mean molar mass $M = \sum_i x_i M_i$; air $\approx 29$~g\,mol$^{-1}$.}
\par\smallskip
\begin{center}
\small
\renewcommand{\arraystretch}{1.12}
\begin{tabular}{@{}lccccl@{}}
\toprule
System (Tab.~5) & $n$ & $T$ & $p$ & $V$ & Law to use\\
\midrule
Weather balloon & fixed & free & free ($= p_{\text{atm}}$) & free & $pV/T = \text{const}$\\
Tyre & fixed & free & free & fixed & $p/T = \text{const}$\\
Cartesian diver & fixed & fixed & free ($\propto$ depth) & free & $pV = \text{const}$\\
Hot-air balloon & free & free & fixed ($= p_{\text{atm}}$) & fixed & $nT = \text{const}$\\
\bottomrule
\end{tabular}
\end{center}

\fbsection{3.C}{Real gases (pp.~20--23)}
\fe[eq.~(15)]{Compressibility factor}{Z = \frac{pV}{nRT} = \frac{pV_m}{RT}}
  {$Z = 1$: ideal gas. $Z < 1$: attractions dominate. $Z > 1$: the volume of the molecules
   dominates. $p \to 0$: $Z \to 1$ (Amagat diagram: $Z$ against $p$).}
\fe[p.~22]{Boyle temperature}{T = T_B\!:\ Z \approx 1 \text{ over a wide range of } p}
  {Below $T_B$, $Z$ first dips below 1; well above, $Z > 1$.}
\fe[eq.~(16)]{Van der Waals}{\left(p + a\,\frac{n^2}{V^2}\right)(V - nb) = nRT}
  {$a$ (Pa\,m$^6$\,mol$^{-2}$): attractions; $b$ (m$^3$\,mol$^{-1}$): co-volume. $a, b \to 0$:
   ideal gas.}
\fe*{Van der Waals, other forms}{\begin{gathered}\left(p + \frac{a}{V_m^2}\right)(V_m - b) = RT\\
  p = \frac{nRT}{V - nb} - a\,\frac{n^2}{V^2}\end{gathered}}
  {Per mole ($V_m = V/n$), and solved for $p$.}
\fe[p.~23]{Internal pressure, co-volume}{\Pi = a\left(\frac{n}{V}\right)^2,\qquad V_{\text{excl}} = nb}
  {$nb$: about four times the volume of the molecules.}
\fe*{Co-volume gas ($a = 0$)}{p\,(V - nb) = nRT}{Repulsions only.}
\fe*{Virial expansion}{Z = 1 + \frac{B(T)}{V_m} + \frac{C(T)}{V_m^2} + \dots}
  {Van der Waals: $B = b - \frac{a}{RT}$, zero at $T_B = \frac{a}{Rb}$.}
\fe*{Van der Waals critical point}{T_c = \frac{8a}{27Rb},\qquad p_c = \frac{a}{27b^2},\qquad V_{m,c} = 3b}
  {Reproduces the $T_c$ column of Tab.~7.}
\par\smallskip
\begin{center}
\small
\renewcommand{\arraystretch}{1.1}
\begin{tabular}{@{}lccccc@{}}
\toprule
Gas (Tab.~7) & He & H$_2$ & N$_2$ & CO$_2$ & H$_2$O\\
\midrule
$a$ (Pa\,m$^6$\,mol$^{-2}$) & $3.46\times10^{-3}$ & $2.45\times10^{-2}$ & 0.137 & 0.364 & 0.554\\
$b$ ($10^{-5}$~m$^3$\,mol$^{-1}$) & 2.38 & 2.67 & 3.87 & 4.27 & 3.05\\
$T_c$ (K) & 5.2 & 33.2 & 126 & 304 & 647\\
\bottomrule
\end{tabular}
\end{center}
\fbnote{\textbf{Other equations of state} (Tab.~6, for reference): Clausius
  $\big(p + \frac{a}{T(V_m + c)^2}\big)(V_m - b) = RT$; Berthelot
  $\big(p + \frac{a}{TV_m^2}\big)(V_m - b) = RT$; Dieterici $p\,(V_m - b)\,e^{a/(V_mRT)} = RT$;
  Soave--Redlich--Kwong $p = \frac{RT}{V_m - b} - \frac{a\,\alpha(T)}{V_m(V_m + b)}$;
  Peng--Robinson $p = \frac{RT}{V_m - b} - \frac{a\,\alpha(T)}{V_m(V_m + b) + b(V_m - b)}$.}

% =====================================================================================
\partheading{Part III\dagnote}{Work, heat and heat engines (beyond Part 1)}
\fbnote{Not in the Part~1 notes: the standard results for ideal-gas cycles. $W$ and $Q$ are
  the energies \textbf{received} by the gas ($> 0$ if received); temperatures in kelvin.}

\fbsection{III.1}{First law and the ideal gas}
\fe*{Work of the pressure forces}{\begin{gathered}\delta W = -p_{\text{ext}}\,\dd V\\
  \text{reversible: } W = -\int_{V_1}^{V_2} p\,\dd V\end{gathered}}
  {$W > 0$ in a compression. Area under the curve in the $(p, V)$ diagram.}
\fe*{First law}{\Delta U = W + Q}
  {Closed system at rest. $U$ is a state function; $W$ and $Q$ depend on the path. Over a
   cycle $\Delta U = 0$.}
\fe*{Internal energy of an ideal gas}{\Delta U = n\,C_{V,m}\,\Delta T}
  {Depends on $T$ only. $C_{V,m} = \frac32 R$ (monatomic), $\frac52 R$ (diatomic near 300~K).}
\fe*{Enthalpy}{H = U + pV,\qquad \Delta H = n\,C_{p,m}\,\Delta T}{At constant pressure: $Q = \Delta H$.}
\fe*{Mayer's relation, $\gamma$}{C_{p,m} - C_{V,m} = R,\qquad \gamma = \frac{C_{p,m}}{C_{V,m}}}
  {$\gamma = 5/3$ (monatomic), $7/5$ (diatomic).}
\par\smallskip
\begin{center}
\small
\renewcommand{\arraystretch}{1.5}
\begin{tabular}{@{}lcccc@{}}
\toprule
Reversible process & Constant & $W$ (received) & $Q$ (received) & $\Delta U$\\
\midrule
Isochoric & $V$ & $0$ & $n\,C_{V,m}\,\Delta T$ & $n\,C_{V,m}\,\Delta T$\\
Isobaric & $p$ & $-p\,\Delta V = -nR\,\Delta T$ & $n\,C_{p,m}\,\Delta T$ & $n\,C_{V,m}\,\Delta T$\\[6pt]
Isothermal & $T$ & $-nRT\ln\dfrac{V_2}{V_1}$ & $nRT\ln\dfrac{V_2}{V_1}$ & $0$\\[9pt]
Adiabatic & $Q = 0$ & $\dfrac{p_2V_2 - p_1V_1}{\gamma - 1}$ & $0$ & $n\,C_{V,m}\,\Delta T$\\
\bottomrule
\end{tabular}
\end{center}
\fe*{Laplace's law}{\begin{gathered}pV^\gamma = \text{const},\qquad TV^{\gamma - 1} = \text{const}\\
  T^\gamma p^{1 - \gamma} = \text{const}\end{gathered}}
  {Reversible adiabatic process of an ideal gas with constant $\gamma$.}

\fbsection{III.2}{Second law and heat engines}
\fe*{Second law (entropy)}{\begin{gathered}\Delta S = S_{\text{exchanged}} + S_{\text{created}},
  \qquad S_{\text{created}} \ge 0\\ S_{\text{exchanged}} = \sum_k \frac{Q_k}{T_k}\end{gathered}}
  {$T_k$: temperature of the source supplying $Q_k$. $S_{\text{created}} = 0$ only for a
   reversible process.}
\fe*{Entropy of an ideal gas}{\Delta S = n\,C_{V,m}\ln\frac{T_2}{T_1} + nR\ln\frac{V_2}{V_1}}{J\,K$^{-1}$.}
\fe*{Heat engine (two sources)}{W + Q_H + Q_C = 0,\qquad \eta = \frac{-W}{Q_H} = 1 + \frac{Q_C}{Q_H}}
  {$Q_H > 0$ taken from the hot source, $Q_C < 0$ given to the cold one, $W < 0$ delivered.}
\fe*{Carnot (maximum) efficiency}{\eta \le \eta_C = 1 - \frac{T_C}{T_H}}
  {Equality for a reversible engine, where $\frac{Q_H}{T_H} + \frac{Q_C}{T_C} = 0$.}
\fe*{Refrigerator and heat pump}{\begin{gathered}e_{\text{fridge}} = \frac{Q_C}{W} \le \frac{T_C}{T_H - T_C}\\
  e_{\text{pump}} = \frac{-Q_H}{W} \le \frac{T_H}{T_H - T_C}\end{gathered}}
  {Coefficients of performance; they can exceed 1.}
\fe*{Carnot cycle of an ideal gas}{\begin{gathered}Q_H = nRT_H\ln\frac{V_2}{V_1},\qquad
  Q_C = -nRT_C\ln\frac{V_3}{V_4}\\ \frac{V_2}{V_1} = \frac{V_3}{V_4}\end{gathered}}
  {Isotherm $T_H$ ($V_1 \to V_2$), adiabat ($\to V_3 = V_{\max}$), isotherm $T_C$
   ($\to V_4$), adiabat back to $V_1 = V_{\min}$. Here $Q_H \neq nRT_H\ln(V_{\max}/V_{\min})$.}
\fe*{Stirling cycle}{\begin{gathered}Q_{\text{in}} = nRT_H\ln\frac{V_{\max}}{V_{\min}}\\
  W = -nR\,(T_H - T_C)\ln\frac{V_{\max}}{V_{\min}}\end{gathered}}
  {Two isotherms and two isochores. With an ideal regenerator, $\eta = 1 - T_C/T_H$. Without
   one, add $n\,C_{V,m}(T_H - T_C)$ to $Q_{\text{in}}$.}

% =====================================================================================
\summaryheading{Summary}{Units of all quantities}
\begin{minipage}[t]{0.485\textwidth}\vspace{0pt}
\small
\renewcommand{\arraystretch}{1.2}
\begin{tabular}[t]{@{}l>{$}l<{$}l@{}}
\toprule
\multicolumn{3}{@{}l}{\bfseries\color{theme}Electricity}\\
\midrule
Charge & q & coulomb, C $=$ A\,s\\
Current & i & ampere, A (SI base unit)\\
Potential, voltage & V,\ u,\ e & volt, V $=$ J\,C$^{-1}$\\
Resistance & R & ohm, $\Omega = $ V\,A$^{-1}$\\
Conductance & G & siemens, S $= \Omega^{-1}$\\
Resistivity & \rho & $\Omega$\,m\\
Conductivity & \sigma & S\,m$^{-1}$\\
Power & p,\ P & watt, W $=$ V\,A\\
Energy & E & joule, J $=$ W\,s\\
Carrier density & n & m$^{-3}$\\
Drift velocity & v & m\,s$^{-1}$\\
Current density & j & A\,m$^{-2}$\\
Efficiency & \eta & no unit\\
\bottomrule
\end{tabular}
\end{minipage}\hfill
\begin{minipage}[t]{0.485\textwidth}\vspace{0pt}
\small
\renewcommand{\arraystretch}{1.2}
\begin{tabular}[t]{@{}l>{$}l<{$}l@{}}
\toprule
\multicolumn{3}{@{}l}{\bfseries\color{theme}Thermodynamics}\\
\midrule
Pressure & p & pascal, Pa $=$ N\,m$^{-2}$\\
Volume & V & m$^3$\\
Temperature & T & kelvin, K\\
Amount of matter & n & mole, mol\\
Molar mass & M & kg\,mol$^{-1}$ (or g\,mol$^{-1}$)\\
Molar volume & V_m & m$^3$\,mol$^{-1}$\\
Density & \rho & kg\,m$^{-3}$\\
Coefficients & \alpha,\ \beta;\ \chi_T & K$^{-1}$; Pa$^{-1}$\\
Gas constant & R & J\,mol$^{-1}$\,K$^{-1}$\\
Van der Waals $a$ & a & Pa\,m$^6$\,mol$^{-2}$\\
Van der Waals $b$ & b & m$^3$\,mol$^{-1}$\\
Energies & W,\ Q,\ U,\ H & joule, J\\
Entropy & S & J\,K$^{-1}$\\
Heat capacities & C_{V,m},\ C_{p,m} & J\,mol$^{-1}$\,K$^{-1}$\\
\bottomrule
\end{tabular}
\end{minipage}

\begin{trapbox}[title={Same letter, different quantity}]
\small
$p$: power (W) in electricity, pressure (Pa) in thermodynamics.\quad
$\rho$: resistivity ($\Omega$\,m) or density (kg\,m$^{-3}$).\quad
$n$: carrier density (m$^{-3}$) or amount of substance (mol).\quad
$\alpha$: expansivity here, never a resistance coefficient.\quad
$E$: energy (J) or electromotive force $E_{\mathrm{Th}}$ (V).
\end{trapbox}

\end{document}
